% GENERATED FILE. Edit canonical lessons, then run npm run generate:books.
% canonical-source-hash: fnv1a64:1dabbd85803b28d6
% canonical-lessons: PT-C381-norma, PT-C381-regulamento, PT-C381-obrigatorio, PT-C381-facultativo, PT-C381-obrigacao

\chapter{As regras}
\label{ch:pt-as-regras}
\hlchaptermodality{381}

\begin{chapteropening}
\textbf{By the end of this chapter:} \emph{I can talk about rules: norms and regulations, what is compulsory or optional, and our obligations.}

Complete the last lesson of chapter 381: I can talk about rules: norms and regulations, what is compulsory or optional, and our obligations.
\end{chapteropening}

\section[a norma]{a norma — the rule, the standard, the norm}

\label{lesson:PT-C381-norma}



\begin{quote}
Before the new one: say the Portuguese for to elect, then the Portuguese for the voter.
\end{quote}

\subsection*{You'll want to know: a norma}
\textbf{a norma} — \textquotedblleft{}the rule, the standard, the norm\textquotedblright{}.

\begin{quote}
A escola tem normas muito claras.
\end{quote}

\begin{quote}
\emph{The school has very clear rules.}
\end{quote}

\emph{Por norma} — as a rule, usually.

\begin{cousinweb}
\emph{norma} is a learned borrowing from Latin \emph{nōrma} \textquotedblleft{}a carpenter's square, a rule, a pattern\textquotedblright{} (Lewis \& Short).
\end{cousinweb}

\begin{grammarlens}[title={What you've built}]
One more word for telling, explaining and describing at B1.
\end{grammarlens}

\subsection*{Your turn}
\begin{itemize}
\raggedright
  \item \emph{Say it:} \emph{a norma}
  \item \emph{Say it:} \emph{a norma}, once more
  \item \emph{Say it:} a sentence of your own with \emph{a norma}
  \item \emph{Recall:} say the Portuguese for to elect, then the Portuguese for the voter, then say \emph{a norma} again
\end{itemize}

\begin{tcolorbox}[breakable,colback=teal!4,colframe=teal!35!black,title={Before you move on}]
What does a norma mean? (\textquotedblleft{}The rule, the standard, the norm\textquotedblright{}.) Say it once more.
\end{tcolorbox}
\section[o regulamento]{o regulamento — the regulations, the rules (of a place)}

\label{lesson:PT-C381-regulamento}



\begin{quote}
Before the new one: say the Portuguese for the voter, then the Portuguese for the rule.
\end{quote}

\subsection*{You'll want to know: o regulamento}
\textbf{o regulamento} — \textquotedblleft{}the regulations, the rules (of a place)\textquotedblright{}.

\begin{quote}
O regulamento do prédio proíbe animais.
\end{quote}

\begin{quote}
\emph{The building's regulations forbid animals.}
\end{quote}

\emph{Cumprir o regulamento} — to follow the regulations.

\begin{cousinweb}
\emph{regulamento} is built on \emph{regular} \textquotedblleft{}to regulate\textquotedblright{}, from Late Latin \emph{rēgulāre}, on \emph{rēgula} \textquotedblleft{}a ruler, a rule\textquotedblright{}, from \emph{regere} \textquotedblleft{}to guide, to rule\textquotedblright{} (Lewis \& Short).
\end{cousinweb}

\begin{grammarlens}[title={What you've built}]
One more word for telling, explaining and describing at B1.
\end{grammarlens}

\subsection*{Your turn}
\begin{itemize}
\raggedright
  \item \emph{Say it:} \emph{o regulamento}
  \item \emph{Say it:} \emph{o regulamento}, once more
  \item \emph{Say it:} a sentence of your own with \emph{o regulamento}
  \item \emph{Recall:} say the Portuguese for the voter, then the Portuguese for the rule, then say \emph{o regulamento} again
\end{itemize}

\begin{tcolorbox}[breakable,colback=teal!4,colframe=teal!35!black,title={Before you move on}]
What does o regulamento mean? (\textquotedblleft{}The regulations, the rules (of a place)\textquotedblright{}.) Say it once more.
\end{tcolorbox}
\section[obrigatório]{obrigatório — compulsory, mandatory}

\label{lesson:PT-C381-obrigatorio}



\begin{quote}
Before the new one: say the Portuguese for the rule, then the Portuguese for the regulations.
\end{quote}

\subsection*{You'll want to know: obrigatório}
\textbf{obrigatório} — \textquotedblleft{}compulsory, mandatory\textquotedblright{}.

\begin{quote}
O uso de capacete é obrigatório.
\end{quote}

\begin{quote}
\emph{Wearing a helmet is compulsory.}
\end{quote}

\emph{É obrigatório} + infinitive: \emph{é obrigatório usar cinto}.

\begin{cousinweb}
\emph{obrigatório} is a learned borrowing from Latin \emph{obligātōrius} \textquotedblleft{}obligatory, binding\textquotedblright{} (EtymDB), from \emph{obligāre} \textquotedblleft{}to bind in obligation\textquotedblright{}, on \emph{ligāre} \textquotedblleft{}to tie\textquotedblright{} (Lewis \& Short).
\end{cousinweb}

\begin{grammarlens}[title={What you've built}]
One more word for telling, explaining and describing at B1.
\end{grammarlens}

\subsection*{Your turn}
\begin{itemize}
\raggedright
  \item \emph{Say it:} \emph{obrigatório}
  \item \emph{Say it:} \emph{obrigatório}, once more
  \item \emph{Say it:} a sentence of your own with \emph{obrigatório}
  \item \emph{Recall:} say the Portuguese for the rule, then the Portuguese for the regulations, then say \emph{obrigatório} again
\end{itemize}

\begin{tcolorbox}[breakable,colback=teal!4,colframe=teal!35!black,title={Before you move on}]
What does obrigatório mean? (\textquotedblleft{}Compulsory, mandatory\textquotedblright{}.) Say it once more.
\end{tcolorbox}
\section[facultativo]{facultativo — optional}

\label{lesson:PT-C381-facultativo}



\begin{quote}
Before the new one: say the Portuguese for the regulations, then the Portuguese for compulsory.
\end{quote}

\subsection*{You'll want to know: facultativo}
\textbf{facultativo} — \textquotedblleft{}optional\textquotedblright{}.

\begin{quote}
A visita ao museu é facultativa.
\end{quote}

\begin{quote}
\emph{The museum visit is optional.}
\end{quote}

The opposite of \emph{obrigatório}.

\begin{cousinweb}
\emph{facultativo} is built on \emph{faculdade} in its sense \textquotedblleft{}power, permission\textquotedblright{}, from Latin \emph{facultās} \textquotedblleft{}capability, means\textquotedblright{} (Lewis \& Short), with the learned suffix \emph{-ivo}.
\end{cousinweb}

\begin{grammarlens}[title={What you've built}]
One more word for telling, explaining and describing at B1.
\end{grammarlens}

\subsection*{Your turn}
\begin{itemize}
\raggedright
  \item \emph{Say it:} \emph{facultativo}
  \item \emph{Say it:} \emph{facultativo}, once more
  \item \emph{Say it:} a sentence of your own with \emph{facultativo}
  \item \emph{Recall:} say the Portuguese for the regulations, then the Portuguese for compulsory, then say \emph{facultativo} again
\end{itemize}

\begin{tcolorbox}[breakable,colback=teal!4,colframe=teal!35!black,title={Before you move on}]
What does facultativo mean? (\textquotedblleft{}Optional\textquotedblright{}.) Say it once more.
\end{tcolorbox}
\section[a obrigação]{a obrigação — the obligation, the duty}

\label{lesson:PT-C381-obrigacao}



\begin{quote}
Before the new one: say the Portuguese for compulsory, then the Portuguese for optional.
\end{quote}

\subsection*{You'll want to know: a obrigação}
\textbf{a obrigação} — \textquotedblleft{}the obligation, the duty\textquotedblright{}.

\begin{quote}
Votar é um direito e uma obrigação.
\end{quote}

\begin{quote}
\emph{Voting is a right and a duty.}
\end{quote}

\emph{Ter a obrigação de} + infinitive — to be obliged to.

\begin{cousinweb}
\emph{obrigação} comes from Latin \emph{obligātiō} \textquotedblleft{}a binding, an obligation\textquotedblright{} (Lewis \& Short), through Old Portuguese \emph{obrigaçon} (EtymDB); the \emph{-r-} for \emph{-l-} shows popular change.
\end{cousinweb}

\begin{grammarlens}[title={What you've built}]
One more word for telling, explaining and describing at B1.
\end{grammarlens}

\subsection*{Your turn}
\begin{itemize}
\raggedright
  \item \emph{Say it:} \emph{a obrigação}
  \item \emph{Say it:} \emph{a obrigação}, once more
  \item \emph{Say it:} a sentence of your own with \emph{a obrigação}
  \item \emph{Recall:} say the Portuguese for compulsory, then the Portuguese for optional, then say \emph{a obrigação} again
\end{itemize}

\begin{tcolorbox}[breakable,colback=teal!4,colframe=teal!35!black,title={Before you move on}]
What does a obrigação mean? (\textquotedblleft{}The obligation, the duty\textquotedblright{}.) Say it once more.
\end{tcolorbox}
